% !TEX root = document.tex
% =====================================================================
% 第6章 数据结构与文件说明 —— 有数据/文件/数据库的软件写，纯算法软件可删
%
% 写什么：
%   1. 输入输出文件格式：字段名、类型、含义，用三线表列成数据字典。
%   2. 内部数据结构：关键类或数据类有哪些字段。
%   3. 目录结构：源码目录树 + 每个目录/文件干什么。
%
% 注意：
%   - 表格用三线表居中，列名和类型要跟代码里的实际字段一致。
%   - 有数据库的软件在这一章写表结构、字段、主外键。
%   - 命令行和界面共用的数据格式只写一次，不要重复。
% =====================================================================

\section{数据结构与文件说明}
\label{sec:data_structure}

\subsection{输入文件格式}
\label{ssec:input_format}
软件接受UTF-8或GB18030编码的CSV文件，第一行为表头，其余每行一条记录。读取时会尝试候选编码列表中的每一种，直到能够完整解码。

\subsection{内部数据结构}
\label{ssec:internal_data}
读取后的数据表示为"表头列表 + 数据行列表"：表头是字符串列表，数据行是字符串列表的列表，二者下标一一对应。统计结果用\lstinline|ColumnSummary|数据类表示，字段含义如表\ref{tab:column_summary}所示。

\begin{table}[H]
\centering
\begin{tabular}{lll}
\toprule
\textbf{字段} & \textbf{类型} & \textbf{含义} \\
\midrule
name & str & 列名 \\
total & int & 该列记录总数 \\
filled & int & 非空单元格数 \\
distinct & int & 不同非空值个数 \\
\bottomrule
\end{tabular}
\caption{ColumnSummary字段说明}
\label{tab:column_summary}
\end{table}

\subsection{输出文件格式}
\label{ssec:output_format}
写出文件统一采用UTF-8带BOM编码的CSV，便于在表格软件中直接打开。报告文件包含\lstinline|column|、\lstinline|total|、\lstinline|filled|、\lstinline|distinct|、\lstinline|missing|五列。

\subsection{源码目录结构}
\label{ssec:dir_structure}
\begin{lstlisting}[language=bash,style=blockstyle]
src/demo_tool/
├── __init__.py      # 版本信息与包导出
├── io.py            # 编码探测与文件读写
├── core.py          # 清洗与筛选规则
├── report.py        # 完整性统计
├── cli.py           # 命令行入口
└── web.py           # 浏览器界面
\end{lstlisting}
各模块职责见第\ref{sec:architecture}章模块说明，函数实现见源程序文档。
